\begin{tabular}{@{}p{2.4cm}p{2.3cm}p{2.2cm}p{7.8cm}@{}}
\hline
Criterio & Valore & Dove & Perche' \\
\hline
$R$ & $\le2\,\mu$m & tutte & validita' del fascio gaussiano: $z_R\simeq2\,\mu$m; oltre $R=2\,\mu$m la particella e' piu' grande della regione ben descritta, la cella e' un'estrapolazione \\
$Pe_r$ & $\ge10$ & acqua e vuoto & la coppia ottica supera di 10 volte il rumore browniano: $Pe_r=\tau/(k_BT_{\rm bath})$, $T_{\rm bath}=310$ K \\
$\Delta T$ & $\le20$ K & acqua & conduzione sferica, Nusselt 1; sopra 20 K convezione e variazione di viscosita' non sono trascurabili \\
$T$ & $\le0{,}7\,T_m$ & vuoto & equilibrio radiativo; per la vaterite $T_m=723$ K (transizione vaterite$\to$calcite), quindi $T\le506$ K \\
$\omega$ & $\le0{,}5\,\omega_{\rm rottura}$ & acqua & il rotore non deve lavorare oltre meta' del limite meccanico; in vuoto $\omega=\tau/\beta_{\rm gas}$ e' formale \\
$P$ & $\le10$ mW & tutte & tetto dichiarato del banco ottico; non e' un limite fisico \\
$\eta_{\rm lim}$ & $\ge10^{-7}$ & tutte & rendimento fisico $\eta_{\rm lim}=\tau\omega_{\rm rottura}/P$; soglia dichiarata di rotore, non una legge \\
trappola & necessaria in vuoto & vuoto & $F_{\rm scat}=P/c$ contro il peso $\rho Vg$; senza supporto la particella non resta nel fascio \\
$t_{\rm spin}$ & $\ge1$ s & vuoto & tempo per arrivare a meta' del limite meccanico a coppia costante; soglia operativa del volano \\
\hline
\end{tabular}
